% !TeX program = pdflatex
% =============================================================================
% Lektionsplan (Diskussionsgrundlage, NICHT das Arbeitsblatt selbst)
% Lorentzkraft Teil 1: Kraft auf stromdurchflossene Leiter
% UZH Praktikum I, Sara Romer Urban, KS im Lee, HS2026
% Klassen: 4fh + 4cdeg
% Induktive Reihenfolge Phaenomen -> makroskopisches Modell -> mikroskopisch
% =============================================================================
\documentclass[11pt,a4paper]{article}

\usepackage[T1]{fontenc}
\usepackage[utf8]{inputenc}
\usepackage[ngerman]{babel}
\usepackage{lmodern}
\usepackage[a4paper,margin=2.2cm]{geometry}
\usepackage{array,booktabs,tabularx,xltabular}
\usepackage{enumitem}
\usepackage{xcolor}
\usepackage{amsmath}
\usepackage{siunitx}
\sisetup{per-mode=fraction,fraction-command=\frac}
\usepackage{url}
\Urlmuskip=0mu plus 1mu
\def\UrlBreaks{\do\/\do-\do_}

\definecolor{titleblue}{HTML}{183A6B}
\definecolor{accentblue}{HTML}{2E6F9E}
\definecolor{groupteal}{HTML}{1B7F72}
\definecolor{darkgray}{HTML}{4A4A4A}
\definecolor{warnred}{HTML}{9E2E2E}

\setlength{\parindent}{0pt}
\setlength{\parskip}{0.5em}

\newcommand{\Methode}[1]{{\small\color{groupteal}\textbf{#1}}}

\usepackage{titlesec}
\titleformat{\section}{\Large\bfseries\color{titleblue}}{\thesection}{0.6em}{}
\titlespacing*{\section}{0pt}{1.2em}{0.5em}

\begin{document}
\sloppy

\begin{center}
{\Huge\bfseries\color{titleblue}Lektionsplan: Lorentzkraft Teil 1}\\[0.2em]
{\Large\bfseries\color{titleblue}Magnetische Kraft auf stromdurchflossene Leiter}\\[0.4em]
{\large Doppellektion (2$\times$45 Min.) -- Diskussionsgrundlage}
\end{center}
\vspace{0.5em}

\noindent\textbf{Klassen:} 4fh (15.09.2026) + 4cdeg (18.09.2026), gemeinsames Material

\noindent\textbf{Lehrperson:} Loris Keller

\vspace{0.3em}
\noindent\textcolor{darkgray}{\textbf{Begriffe in diesem Dokument:} die
makroskopische Kraft auf einen stromdurchflossenen Leiter im Feld
($F=IlB$) heisst \textbf{magnetische Kraft}; die Kraft auf eine einzelne
bewegte Ladung ($F=qvB$) heisst \textbf{Lorentzkraft}. Die magnetische
Kraft ist die Summe der Lorentzkraft über alle Ladungsträger im Leiter --
beide Begriffe werden im Folgenden konsequent unterschieden.}

\vspace{0.3em}
\noindent\textcolor{darkgray}{Dieses Dokument ist \textbf{keine} Reinschrift eines Schülerdokuments, sondern die
Planungsgrundlage dafür: Lernziele, Aufbau der Lektion und die gewählte
Didaktik-Methode pro Phase. Nach Absprache werden daraus lokal in VS Code
\textbf{drei separate Dokumente} erstellt: (1) ein
\textbf{Experimente-Dokument} (Protokoll/Anleitung für Leiterschaukel,
zwei Drähte und Fadenstrahlrohr), (2) ein \textbf{Arbeitsblatt}
(Theorie/Erarbeitung für die Lektion selbst) und (3) ein
\textbf{Aufgabenblatt} (zusätzliche Übungsaufgaben) -- siehe Abschnitt 4.
Direkte Fortsetzung von Magnetfelder (4fh 01.09.2026) -- die Feldformel
des geraden Leiters wird hier direkt weiterverwendet. Teil 2 (spätere
Doppellektion) behandelt die Lorentzkraft auf einzelne \emph{freie}
bewegte Ladungen (Kreisbahn, Fadenstrahlrohr, Anwendungen).}

% =========================================================================
\section*{1. Abgrenzung und Lernziele}

\textbf{Im Fokus:} induktiver Aufbau vom Phänomen zum Modell -- zuerst
Beobachtung an zwei Experimenten (Leiterschaukel, zwei parallele Leiter),
dann die \textbf{magnetische Kraft} auf den Leiter als makroskopisches
Modell ($F=I\cdot l\cdot B$, Drei-Finger-Regel), dann die Herleitung der
eigentlichen \textbf{Lorentzkraft} $F=qvB$ (Kraft auf eine \emph{einzelne}
bewegte Ladung), die zeigt, warum die magnetische Kraft gilt und warum
sich die zwei Leiter je nach Stromrichtung anziehen oder abstossen.
Begrifflich wichtig: $F=qvB$ \emph{ist} die Lorentzkraft, $F=IlB$ heisst
magnetische Kraft und ist nur ihre Summe für viele Ladungsträger in einem
Leiter. Direkt nach der Herleitung der Lorentzkraft wird zuerst der Aufbau
des Fadenstrahlrohrs erklärt (was mit den Elektronen passiert), bevor ein
Predict-Observe-Explain-Zyklus dieselbe Lorentzkraft in ihrer reinsten
Form an einer freien Ladung zeigt und die SuS die Kreisbewegung der
Elektronen als Zentripetalkraft-Wirkung der Lorentzkraft erklären --
bevor danach zu den zwei Leitern
zurückgekehrt wird, um dieselbe mikroskopische Perspektive auf ein
zweites Phänomen (Anziehung/Abstossung) anzuwenden. Am Ende der Lektion
mündet das Ganze in einen kurzen Ausblick auf Teil 2, das die Lorentzkraft
auf freie Ladungen quantitativ vertieft.

\textbf{Bereits bekannt (nur Repetition, keine Neueinführung):} das
Magnetfeld eines geraden, stromdurchflossenen Leiters
($B=\frac{\mu_0}{2\pi}\cdot\frac{I}{r}$, Magnetfelder-Doppellektion),
Feldlinien, $\odot$/$\otimes$-Konvention; zwei bereits bekannte
rechte-Hand-Regeln (gerader Leiter, Spule) -- wichtiger Bezugspunkt, weil
mit der Drei-Finger-Regel eine \emph{dritte}, andersartige Handregel hinzukommt,
diesmal zusätzlich in zwei Varianten (rechts/positiv, links/negativ);
Zentripetalkraft (Mechanik) -- wird bereits hier für die qualitative
Erklärung der Kreisbahn im Fadenstrahlrohr gebraucht (LZ6) und in Teil 2
quantitativ vertieft ($r=mv/qB$); das
\textbf{Kreuzprodukt inkl.\ Sinus-Betragsformel} ($|\vec a\times\vec
b|=ab\sin\alpha$, aus Mathematik) -- neu bestätigtes Vorwissen, das den
allgemeinen Winkel-Fall mathematisch zugänglich macht.

\textbf{Bewusst nicht Teil dieser Einheit (Kernstoff):} die vollständige
Vektorform $\vec F=I\cdot\vec l\times\vec B$ bzw.\ $\vec F=q\cdot\vec
v\times\vec B$ ist kein eigenständiges Lernziel -- im Zentrum steht
weiterhin der senkrechte Fall ($\vec I\perp\vec B$ bzw.\ $\vec v\perp\vec
B$) mit der Drei-Finger-Regel, da die Experimente selbst senkrecht aufgebaut
sind; die allgemeine Formel $F=IlB\sin\alpha$ wird trotzdem kurz in der
Theorie erwähnt, da das Kreuzprodukt bereits bekannt ist, dazu Aufgaben
mit beliebigem Winkel im Aufgabenblatt (siehe Abschnitt 4/5). Ebenfalls
nicht Teil: die \emph{quantitative} Lorentzkraft auf einzelne \emph{freie}
bewegte Ladungen (Bahnradius $r=mv/qB$, sechs Anwendungen -- Teil 2),
Induktion, Elektromotor-Funktionsprinzip im Detail. Die \emph{qualitative}
Erklärung der Kreisbahn im Fadenstrahlrohr (LZ6) ist davon ausgenommen und
findet bereits hier statt.

\begin{itemize}[leftmargin=1.3cm,itemsep=0.3em]
  \item[LZ1] die Beobachtungen an der Leiterschaukel und an zwei
  parallelen Leitern beschreiben, ohne diese bereits zu erklären.
  \item[LZ2] die Formel $F=I\cdot l\cdot B$ für die \textbf{magnetische
  Kraft} auf einen stromdurchflossenen Leiter im (senkrechten) Magnetfeld
  anwenden.
  \item[LZ3] die Richtung der magnetischen Kraft mit der Drei-Finger-Regel
  (Daumen = Strom/Ursache, Zeigefinger = Feld/Vermittlung,
  Mittelfinger = Kraft/Wirkung; rechte Hand, da die Stromrichtung per
  Konvention die Bewegungsrichtung positiver Ladungsträger ist) bestimmen
  und explizit von den beiden bereits bekannten rechte-Hand-Regeln
  unterscheiden. Ab LZ5 kommt dieselbe Regel auch als
  \textbf{linke-Hand-Variante für negative Ladungsträger} zum Einsatz.
  \item[LZ4] die eigene Beobachtung am Leiterschaukel-Experiment (LZ1)
  mithilfe von LZ2 und LZ3 (magnetische Kraft, Drei-Finger-Regel) selbst erklären.
  \item[LZ5] die \textbf{Lorentzkraft} $F=qvB$ als Kraft auf eine
  \emph{einzelne} bewegte Ladung herleiten und ihre \textbf{Richtung
  ebenfalls mit der Drei-Finger-Regel bestimmen} -- rechte Hand für
  \emph{positive}, linke Hand für \emph{negative} Ladungen (Daumen =
  Geschwindigkeit $v$ der Ladung statt Stromrichtung), sodass sich die
  Richtung für beide Vorzeichen direkt ablesen lässt, ohne einen
  zusätzlichen Umkehr-Schritt; zeigen, dass die Summe über $Q=It$ wieder die
  magnetische Kraft $F=IlB$ ergibt -- die magnetische Kraft ist die
  makroskopische Konsequenz der Lorentzkraft für einen Leiter, nicht die
  Lorentzkraft selbst.
  \item[LZ6] die Kreisbewegung der Elektronen im Fadenstrahlrohr
  \textbf{erklären}, indem die Lorentzkraft als \textbf{Zentripetalkraft}
  gedeutet wird (stets senkrecht zu $v$ $\Rightarrow$ nur Richtungs-, kein
  Betragswechsel von $v$ $\Rightarrow$ Kreisbahn); die Krümmungsrichtung
  mit der linken-Hand-Variante der Drei-Finger-Regel bestimmen (Elektronen
  negativ geladen) -- im Rahmen eines Predict-Observe-Explain-Zyklus im
  Plenum.
  \item[LZ7] die Kraft zwischen zwei parallelen stromdurchflossenen
  Leitern herleiten (Feldformel des geraden Leiters kombiniert mit der
  magnetischen Kraft $F=I\cdot l\cdot B$: $\frac{F}{l}=\frac{\mu_0}{2\pi}
  \cdot\frac{I_1\cdot I_2}{r}$).
  \item[LZ8] erklären, wieso gleichgerichtete Ströme sich anziehen und
  entgegengesetzt gerichtete Ströme sich abstossen -- mithilfe von LZ5
  (Drei-Finger-Regel an einer einzelnen Ladung im Feld des Nachbarleiters).
\end{itemize}

% =========================================================================
\section*{2. Aufbau der Doppellektion}

\begin{xltabular}{\textwidth}{@{}p{2.6cm}X p{1.2cm}@{}}
\toprule
\textbf{Phase} & \textbf{Inhalt \& Methode} & \textbf{Zeit} \\
\midrule
\endhead

Einstieg &
\textbf{Advance Organizer: Rückblick + Alltagsbeispiele + Roadmap.}
Rückblick: bei Magnetfelder gelernt, dass bewegte Ladungen (Ströme) ein
Magnetfeld erzeugen -- neue Frage: was passiert, wenn eine bewegte Ladung
sich selbst in einem Magnetfeld befindet? Explizit klären: gemeint ist ein
\textbf{externes} Magnetfeld (von aussen, z.\,B.\ von einem Magneten oder
einem anderen Leiter), \emph{nicht} das Feld, das die Ladung selbst
erzeugt. Danach zwei Bilder als Aufhänger: Lautsprecher/Elektromotor
(stromdurchflossene Spule bewegt sich im Feld -- magnetische Kraft) und
Polarlicht (geladene Teilchen aus dem Sonnenwind werden vom Erdmagnetfeld
auf gekrümmte Bahnen gelenkt -- Lorentzkraft). Kurze Fragen dazu,
\emph{ohne} aufzulösen, dann Roadmap als Tafelskizze (3 Stationen:
magnetische Kraft $\to$ Lorentzkraft $\to$ freie Ladung in Teil 2), die
während der Lektion sichtbar bleibt -- Details und vollständiger Text
siehe Abschnitt 2.1. \Methode{A1 Advance Organizer} -- Lehrervortrag mit
kurzem Unterrichtsgespräch, Plenum. &
3' \\[0.4em]

Aktivierung &
\textbf{Leiterschaukel: reine Beobachtung.} Leiterschaukel im Magnetfeld
(umgebauter Magnetfelder-Baukasten, Spule entfernt), Stromkreis umpolbar.
Strom anschliessen, Richtung der Auslenkung bei unterschiedlicher
Strom-/Feldpolung beobachten und notieren -- bewusst nur Beobachtung,
noch keine Erklärung (LZ1). Sozialform unbestätigt (siehe Abschnitt 5). &
6' \\[0.4em]

Aktivierung &
\textbf{Zwei parallele Leiter: reine Beobachtung.} Zwei stromdurch-
flossene Drähte, je nach Stromrichtung Anziehung oder Abstossung -- auch
hier bewusst nur Beobachtung, noch keine Erklärung (LZ1). Sozialform
unbestätigt (siehe Abschnitt 5). &
5' \\[0.4em]

Theorie &
\textbf{Magnetische Kraft: Formel und Drei-Finger-Regel.} $F=I\cdot l\cdot B$ für
den senkrechten Fall als magnetische Kraft auf den Leiter einführen
(LZ2). Richtung über die Drei-Finger-Regel (rechte Hand, drei paarweise
senkrechte Finger) -- explizit gegen die beiden bereits bekannten
rechte-Hand-Regeln abgrenzen (LZ3); kurzer Hinweis, dass ab der
Lorentzkraft dieselbe Regel auch mit der linken Hand für negative
Ladungsträger verwendet wird. Kurze Erwähnung des allgemeinen
Winkel-Falls ($F=IlB\sin\alpha$, Kreuzprodukt bereits aus der Mathematik
bekannt) -- hier aber stets $\alpha=90^\circ$. Lehrervortrag, geleitet,
Tafelskizze. &
10' \\[0.4em]

Rückkehr &
\textbf{Leiterschaukel: eigene Beobachtung erklären.} Mit der neuen
Formel und der Drei-Finger-Regel die eigene Beobachtung aus der
Leiterschaukel-Aktivierung selbst erklären (LZ4). \Methode{C3 Think-Pair-Share} -- zuerst allein
anwenden, dann zu zweit vergleichen, danach Plenum. &
5' \\[0.4em]

Übung &
\textbf{Anwendung der magnetischen Kraft und Drei-Finger-Regel.} Kurze Aufgaben:
Kraft berechnen, Richtung mit der Drei-Finger-Regel bestimmen. Individuelles
Üben. &
6' \\[0.4em]

Vertiefung &
\textbf{Herleitung der Lorentzkraft.} Warum wirkt überhaupt eine Kraft auf
den ganzen Leiter? Einzelne Ladungsträger im Feld erfahren je die
Lorentzkraft $F=qvB$ (dieselbe Drei-Finger-Regel, jetzt auf eine Ladung
angewendet -- rechte Hand für positive, linke Hand für negative Ladungen);
mit $Q=It$ und $v=l/t$ zeigen, dass die Summe der
Einzelkräfte wieder die magnetische Kraft $F=IlB$ ergibt (LZ5) -- $F=qvB$
\emph{ist} die Lorentzkraft, $F=IlB$ die magnetische Kraft -- nur ihre
Summe für den Leiter. Gemeinsam entwickelte Kurzherleitung im Plenum. &
12' \\[0.4em]

Theorie &
\textbf{Wie funktioniert das Fadenstrahlrohr?} Aufbau erklären, \emph{ohne}
die Kreisbahn vorwegzunehmen: die Elektronen werden durch eine Spannung
beschleunigt, treten dann in ein (näherungsweise homogenes) Magnetfeld
ein und werden durch das Restgas im Rohr als leuchtende Bahn sichtbar
gemacht. Lehrervortrag, Tafelskizze, Plenum. &
4' \\[0.4em]

POE &
\textbf{Predict-Observe-Explain: Fadenstrahlrohr.} \emph{Predict:} die SuS sagen mithilfe
von $F=qvB$ (Kraft stets senkrecht zu $v$) und der aus der Mechanik
bekannten Zentripetalkraft eine Kreisbahn voraus. \emph{Observe:} das Feld
wird eingeschaltet, die Elektronen laufen sichtbar auf einer Kreisbahn --
dieselbe Lorentzkraft, jetzt an einer \emph{freien} Ladung ohne
Leiterwand, in ihrer reinsten Form gezeigt. \emph{Explain:} die SuS
erklären im Plenum, wie sich die Elektronen verhalten -- die Lorentzkraft
wirkt als Zentripetalkraft, die Krümmungsrichtung folgt aus der
linken-Hand-Variante der Drei-Finger-Regel (Elektronen negativ geladen,
LZ6). \Methode{M4 Predict-Observe-Explain} -- Lehrer-Demonstration,
Plenum. &
8' \\[0.4em]

Theorie &
\textbf{Rückkehr zu den zwei Leitern: Herleitung.} Feld des ersten
Leiters am Ort des zweiten
($B_1(r)=\frac{\mu_0}{2\pi}\cdot\frac{I_1}{r}$, bereits bekannt aus der
Magnetfelder-Doppellektion) einsetzen in
$F=I_2\cdot l\cdot B_1$ ergibt
$\frac{F}{l}=\frac{\mu_0}{2\pi}\cdot\frac{I_1\cdot I_2}{r}$ (LZ7).
Gemeinsame Herleitung an der Tafel, ggf. Teilschritte in Partnerarbeit. &
10' \\[0.4em]

Theorie &
\textbf{Anziehung vs. Abstossung -- mikroskopisch erklärt.} Drei-Finger-Regel an
einer einzelnen Ladung im Feld des jeweils anderen Leiters anwenden:
gleiche Stromrichtung $\Rightarrow$ Anziehung, entgegengesetzte Richtung
$\Rightarrow$ Abstossung (LZ8) -- dieselbe mikroskopische Perspektive wie
beim Fadenstrahlrohr und in der Vertiefungsphase, jetzt auf das zweite
Phänomen angewendet; explizit gegen die Intuition \glqq gleich stösst
ab\grqq{} (Pole, Ladungen) abgrenzen. Lehrervortrag, geleitet. &
6' \\[0.4em]

Ausblick &
\textbf{Historischer Kontext (optional).} Dieser Aufbau diente früher
direkt zur Definition der Einheit Ampere (vor der SI-Neudefinition 2019)
-- Quelle/Details noch zu verifizieren. Als Erstes streichbar bei
Zeitdruck. &
3' \\[0.4em]

Abschluss &
\textbf{Zusammenfassung und Ausblick auf Teil 2/Aufgabenblatt.} Kurze
Zusammenfassung der Lektion, kurzer Ausblick, dass Teil 2 die Lorentzkraft
auf freie Ladungen (Kreisbahn, sechs Anwendungen) systematisch vertieft,
Hinweis auf das separate Aufgabenblatt für weitere Übung. Plenum. &
3' \\

\bottomrule
\end{xltabular}

\textcolor{darkgray}{\small Summe: 81 Min. -- 9 Min. Puffer innerhalb der
90 Min. der Doppellektion -- siehe Zeitbudget in Abschnitt 5.}

\subsection*{2.1 Ausgestaltungsvorschlag: Advance Organizer}

Ziel: in 3 Minuten Interesse wecken und die Roadmap der Lektion zeigen,
\emph{ohne} die Experimente inhaltlich vorwegzunehmen -- mit zwei
konkreten Alltagsbeispielen statt einer abstrakten Ankündigung, eines pro
Kraft-Typ dieser Lektion.

\textbf{Rückblick als Einstieg (wichtig).} Bei Magnetfelder haben die SuS
gelernt: bewegte Ladungen (Ströme) erzeugen ein Magnetfeld. Die neue Frage
dieser Lektion dreht das um: Was passiert, wenn eine bewegte Ladung sich
selbst in einem Magnetfeld befindet? Zentral dabei: gemeint ist ein
\textbf{externes} Magnetfeld -- eines, das von aussen kommt (z.\,B.\ von
einem Permanentmagneten oder einem anderen Leiter) -- \emph{nicht} das
Magnetfeld, das die bewegte Ladung selbst erzeugt. Diese Unterscheidung
sollte explizit angesprochen werden, da sie sonst leicht verwechselt wird
(die SuS kennen aus Magnetfelder nur die Richtung \glqq Ladung/Strom
$\to$ Feld\grqq{}, jetzt kommt \glqq Feld $\to$ Kraft auf die Ladung\grqq{}
dazu).

\textbf{Beispiel 1 (magnetische Kraft, $F=IlB$): Lautsprecher oder
Elektromotor.} In jedem Lautsprecher sitzt eine stromdurchflossene Spule
in einem starken Magnetfeld; fliesst Strom, bewegt sich die Spule und
damit die Membran -- genau das erzeugt den Ton. Derselbe Mechanismus
steckt im Elektromotor (Rasierapparat, Ventilator, Fensterheber,
elektrische Zahnbürste). Frage an die Klasse: \glqq Was hat ein
Lautsprecher mit einem stromdurchflossenen Draht in einem Magnetfeld zu
tun?\grqq{}

\textbf{Beispiel 2 (Lorentzkraft, $F=qvB$): Polarlicht (Nordlicht/Aurora).}
Geladene Teilchen aus dem Sonnenwind treffen auf das Erdmagnetfeld und
werden davon auf spiralförmige Bahnen zu den Polen abgelenkt -- dort
leuchten sie beim Auftreffen auf die Atmosphäre. Hier wirkt dieselbe Art
Kraft nicht auf einen ganzen Draht, sondern auf \emph{einzelne} freie
geladene Teilchen. Frage an die Klasse: \glqq Warum bewegen sich diese
Teilchen nicht geradeaus, sondern auf einer gekrümmten Bahn?\grqq{}

\textbf{Konkreter Ablauf (ca.\ 3'):}
\begin{itemize}[leftmargin=1.2cm,itemsep=0.2em]
  \item (ca.\ 40'') Rückblick + neue Frage stellen: bewegte Ladungen
  erzeugen ein Magnetfeld (bekannt) -- was passiert umgekehrt, wenn eine
  bewegte Ladung sich in einem \textbf{externen} Magnetfeld befindet?
  Explizit klarstellen, dass \emph{nicht} das selbst erzeugte Feld gemeint
  ist.
  \item (ca.\ 40'') Beide Bilder zeigen (Folie/Ausdruck/Skizze), die
  beiden Fragen stellen, kurz im Raum stehen lassen -- \emph{ohne}
  aufzulösen. Ziel ist Neugier, keine Antwort.
  \item (ca.\ 40'') Überleitung: \glqq Beide Phänomene haben denselben
  Ursprung: eine Kraft, die externe Magnetfelder auf Ströme bzw.\ auf
  bewegte Ladungen ausüben. Genau die untersuchen wir jetzt zwei
  Doppellektionen lang.\grqq{}
  \item (ca.\ 60'') Roadmap als knappe Tafelskizze (3 Stationen), die
  während der ganzen Lektion sichtbar bleiben kann: \glqq Zuerst schauen
  wir uns die Kraft auf den ganzen Leiter an -- die magnetische Kraft, wie
  beim Lautsprecher/Elektromotor. Danach fragen wir, was auf der Ebene der
  einzelnen Ladung passiert -- das ist die Lorentzkraft, wie beim
  Polarlicht. Und in der nächsten Doppellektion schauen wir, was passiert,
  wenn eine Ladung nicht mehr im Leiter gebunden, sondern ganz frei ist --
  so wie beim Polarlicht oder in einem Teilchenbeschleuniger.\grqq{}
\end{itemize}

\textbf{Weitere mögliche Alltagsbeispiele} (Ersatz/Ergänzung): ein alter
Röhrenbildschirm/Röhrenfernseher (Elektronen werden durch Magnetspulen
abgelenkt -- historischer Vorläufer des Fadenstrahlrohrs, ohne das hier
schon zu verraten); ein Massenspektrometer oder ein Teilchenbeschleuniger
(z.\,B.\ CERN) als \glqq wozu braucht man das in der Forschung?\grqq{}-Ausblick.

\subsection*{2.2 Kerninhalte und Leitfragen je Phase}

Kurzüberblick über die wichtigste inhaltliche Botschaft und die Leitfrage
jeder Phase aus Abschnitt 2 -- als schneller Fahrplan, auch für die
spätere Arbeitsblatt-Erstellung.

\begin{itemize}[leftmargin=1.3cm,itemsep=0.3em]
  \item \textbf{Einstieg (Advance Organizer):} bewegte Ladungen erzeugen
  ein Magnetfeld (bekannt) -- was passiert umgekehrt, wenn eine bewegte
  Ladung sich in einem \emph{externen} Magnetfeld befindet? Dieselbe Kraft
  steckt in Lautsprecher/Elektromotor und in Polarlichtern.
  \emph{Leitfrage:} Was passiert, wenn eine bewegte Ladung sich in einem
  externen (nicht selbst erzeugten) Magnetfeld befindet?
  \item \textbf{Aktivierung Leiterschaukel:} ein stromdurchflossener
  Leiter im Feld wird ausgelenkt, je nach Strom-/Feldrichtung anders.
  \emph{Leitfrage:} Was passiert mit einem stromdurchflossenen Leiter im
  Magnetfeld?
  \item \textbf{Aktivierung zwei Leiter:} zwei stromdurchflossene Leiter
  ziehen sich an oder stossen sich ab. \emph{Leitfrage:} Warum ziehen sich
  zwei stromdurchflossene Leiter manchmal an und stossen sich manchmal ab?
  \item \textbf{Theorie Magnetische Kraft:} $F=IlB$ als Modell,
  Drei-Finger-Regel (rechte Hand) für die Richtung, allgemeiner
  Winkel-Fall kurz erwähnt. \emph{Leitfrage:} Wie berechnet man die Kraft
  auf den Leiter, und wie findet man ihre Richtung?
  \item \textbf{Rückkehr Leiterschaukel:} die eigene Beobachtung wird mit
  Formel und Regel selbst erklärt. \emph{Leitfrage:} Kann ich meine eigene
  Beobachtung jetzt selbst erklären?
  \item \textbf{Übung:} Anwendung von Formel und Drei-Finger-Regel an
  Beispielen. \emph{Leitfrage:} Wie wende ich Formel und Regel sicher an?
  \item \textbf{Vertiefung Herleitung Lorentzkraft:} $F=IlB$ ist die Summe
  der Lorentzkraft $F=qvB$ über alle einzelnen Ladungsträger.
  \emph{Leitfrage:} Warum wirkt die Kraft überhaupt auf den ganzen Leiter
  -- was passiert dabei auf der Ebene der einzelnen Ladungen?
  \item \textbf{Theorie Fadenstrahlrohr-Aufbau:} Funktionsweise
  (Beschleunigung, Eintritt ins Feld, Sichtbarmachung durch Restgas) --
  ohne das Ergebnis vorwegzunehmen. \emph{Leitfrage:} Wie kann man die
  Bewegung einer einzelnen freien Ladung überhaupt sichtbar machen?
  \item \textbf{POE Fadenstrahlrohr:} die Lorentzkraft wirkt als
  Zentripetalkraft und erzeugt eine Kreisbahn; Richtung über die linke
  Hand (negative Ladung). \emph{Leitfrage:} Wie bewegt sich eine freie
  Ladung, auf die die Lorentzkraft wirkt?
  \item \textbf{Theorie Rückkehr zwei Leiter (Herleitung):} das
  Kraftgesetz zwischen zwei parallelen Leitern ergibt sich aus Feldformel
  und magnetischer Kraft kombiniert. \emph{Leitfrage:} Wie hängt die Kraft
  zwischen zwei Leitern von Abstand und Stromstärken ab?
  \item \textbf{Theorie Anziehung vs.\ Abstossung:} mikroskopische
  Erklärung über die Lorentzkraft auf eine Ladung im Feld des
  Nachbarleiters. \emph{Leitfrage:} Warum genau zieht gleiche
  Stromrichtung an, entgegengesetzte stösst ab?
  \item \textbf{Ausblick historischer Kontext (optional):} dieser Aufbau
  diente früher direkt zur Definition der Einheit Ampere. \emph{Leitfrage:}
  Wie wurde die Stromstärke-Einheit eigentlich früher definiert?
  \item \textbf{Abschluss:} Zusammenfassung, Ausblick auf die quantitative
  Vertiefung in Teil 2. \emph{Leitfrage:} Was passiert mit derselben
  Kraft, wenn die Ladung ganz frei ist?
\end{itemize}

% =========================================================================
\section*{3. Begründung der Methodenwahl}

Beide Experimente sind bewusst als \glqq Beobachten zuerst, Erklären
später\grqq{} angelegt, nicht als vollständiger Predict-Observe-Explain-
Zyklus (M4): Vor der Theorie fehlt den SuS die Grundlage für eine
begründete Vorhersage -- eine Vorhersage ohne Konzept würde nur zu
zufälligem Raten führen, nicht zu echtem Vorwissen-Konflikt. Die
Rückkehr zum Leiterschaukel-Experiment nach der Theorie (C3 Think-Pair-
Share) schliesst den Bogen aber explizit: die SuS wenden die neue Regel
aktiv auf ihre eigene, bereits gemachte Beobachtung an, statt nur ein
neues Beispiel zu rechnen.

Die Herleitung der Lorentzkraft folgt bewusst \emph{nach} der magnetischen
Kraft als makroskopischem Modell, nicht davor: Die SuS kennen zu diesem
Zeitpunkt bereits $F=IlB$ und die Drei-Finger-Regel und können die Lorentzkraft
$F=qvB$ deshalb als vertiefende Begründung des bereits Verstandenen
einordnen,
statt als isolierte neue Formel -- auch wenn $F=qvB$ physikalisch die
grundlegendere Grösse ist. Direkt danach wird dieselbe Lorentzkraft am
Fadenstrahlrohr an einer freien Ladung gezeigt (zweite Anwendung), bevor
die mikroskopische Perspektive ein drittes Mal auf die zwei Leiter
angewendet wird (Anziehung/Abstossung) -- die Wiederholung an mehreren
Beispielen in kurzer Folge soll die neue Denkweise festigen und zugleich
zeigen, wie weit dieselbe Kraft trägt: vom Leiter bis zur freien Ladung,
bevor die Lektion wieder zum Leiter zurückkehrt.

Das Fadenstrahlrohr wird bewusst in zwei getrennten Phasen behandelt statt
in einer einzigen Demonstration: zuerst der Aufbau (was mit den
Elektronen passiert), \emph{ohne} die Kreisbahn vorwegzunehmen, danach ein
vollständiger Predict-Observe-Explain-Zyklus (M4). Diese Trennung stellt
sicher, dass die SuS beim \emph{Predict} tatsächlich etwas vorhersagen
(die Kraft steht senkrecht zu $v$, ist also analog zur bereits bekannten
Zentripetalkraft) und nicht nur eine unbekannte Apparatur bestaunen -- die
Vorhersage braucht das Wissen über den Aufbau, aber nicht das Ergebnis.
Damit wird die qualitative Kreisbahn-Erklärung bereits hier abschliessend
geleistet; Teil 2 kann direkt mit der quantitativen Vertiefung
($r=mv/qB$) einsteigen, statt dieselbe Frage nochmals zu stellen (siehe
Abschnitt 5).

Die Herleitung der Kraft zwischen zwei Leitern (LZ7) verknüpft bewusst
zwei bereits an anderer Stelle eingeführte Formeln (Feld des geraden
Leiters, Kraft auf einen Leiter im Feld) statt eine neue Formel direkt
vorzugeben -- die SuS sehen dadurch, dass die neue Grösse keine
eigenständige Regel ist, sondern eine direkte Konsequenz aus bereits
Bekanntem.

% =========================================================================
\section*{4. Anschlussdokumente: Experimente, Arbeitsblatt, Aufgabenblatt}

Aus diesem Lektionsplan entstehen drei separate Dokumente:

\begin{itemize}[leftmargin=1.2cm,itemsep=0.4em]
  \item \textbf{Experimente-Dokument:} Protokoll/Anleitung für die drei
  Experimente dieser Lektion (Leiterschaukel, zwei parallele Leiter,
  Fadenstrahlrohr) -- Aufbau, Durchführung, was die SuS jeweils
  beobachten und notieren sollen.
  \item \textbf{Arbeitsblatt:} Theorie- und Erarbeitungsdokument für die
  Lektion selbst (analog zu den bisherigen Themen), entlang der in
  Abschnitt 2 beschriebenen Phasen.
  \item \textbf{Aufgabenblatt:} zusätzliche Übungsaufgaben zur
  magnetischen Kraft ($F=I\cdot l\cdot B$), zur Richtungsbestimmung
  (rechte/linke Hand je nach Vorzeichen), zur
  Lorentzkraft, zum allgemeinen Winkel-Fall ($F=IlB\sin\alpha$) und zur
  Kraft zwischen zwei parallelen Leitern.
\end{itemize}

% =========================================================================
\section*{5. Offene Punkte zur Besprechung}

\begin{itemize}[leftmargin=1.2cm,itemsep=0.4em]
  \item \textcolor{warnred}{\textbf{Konsequenz für Teil 2:}} Vor dem
  Fadenstrahlrohr wird zuerst der Aufbau erklärt und danach ein
  vollständiger Predict-Observe-Explain-Zyklus mit demselben Experiment
  durchgeführt (LZ6). Damit ist die qualitative Frage \glqq bewegt sich
  die Ladung im Kreis?\grqq{} nach Teil 1 bereits beantwortet und gezeigt.
  Der für Teil 2 geplante Einstieg
  (\texttt{diskussion-teil1-teil2-inhalte.md}) darf diese Frage deshalb
  nicht nochmals als Überraschung anlegen, sondern sollte direkt auf die
  \textbf{quantitative} Vertiefung fokussieren: Herleitung des Bahnradius
  ($r=mv/qB$) und seine Abhängigkeit von Spannung, Feldstärke und
  spezifischer Ladung, sowie die sechs Anwendungen. Bitte
  \texttt{diskussion-teil1-teil2-inhalte.md} entsprechend anpassen.
  \item \textbf{Sozialform der beiden Experimente:} Sind Leiterschaukel
  und/oder Zwei-Leiter-Aufbau als Schülerexperiment (mehrere Sets,
  Partner-/Gruppenarbeit, wie beim Magnetfelder-Schülerexperiment) oder
  als Lehrer-Demonstration (ein Aufbau, ähnlich der Ørsted-Demo bei
  Magnetfelder) gedacht? Beeinflusst Zeitbudget und Methodenwahl direkt.
  \item \textbf{Reihenfolge der beiden Beobachtungen:} Sollen beide
  Beobachtungen nacheinander als zwei Stationen laufen (wie hier
  geplant, gemeinsame Theorie danach für beide), oder soll zuerst die
  Leiterschaukel komplett durchlaufen werden (Beobachtung $\to$ Theorie
  $\to$ mikroskopisch) und erst danach separat die zwei Leiter?
  \item \textbf{PhysikLibre-Referenz:} Kapitel 13.5 \glqq Leiter in
  Magnetfeldern\grqq{} liegt bereits als Referenz-PDF im Repo vor --
  genaue Seitenzahlen/Abschnitte innerhalb des Kapitels wurden aber noch
  nicht im Detail gegengeprüft.
  \item \textbf{Vorwissen-Material:} die Datei
  \texttt{Werkstatt\_Magnetismus.docx} (in
  \texttt{elektrodynamik/Vorwissen/Magnetismus/}) liegt im Repo vor,
  wurde aber inhaltlich noch nicht konsultiert -- falls Sara Romer dort
  bereits
  Apparatur-Fotos, verwendete Formulierungen oder einen bevorzugten
  Schwierigkeitsgrad festgehalten hat, sollte das vor der Arbeitsblatt-
  Erstellung noch eingearbeitet werden.
  \item \textbf{Historischer Ampere-Exkurs:} Nice-to-have, nicht
  zeitkritisch -- als Erstes streichbar, falls die Lektion zu voll wird.
  \item \textbf{Zeitbudget:} Elf Phasen mit Fachinhalt (Advance Organizer,
  zwei Beobachtungen, makroskopisches Modell, Rückkehr, Übung, Herleitung
  der Lorentzkraft, Fadenstrahlrohr-Aufbau, Fadenstrahlrohr-POE, Herleitung
  Kraft zwischen Leitern, mikroskopische Erklärung Anziehung/Abstossung)
  ergeben bei 9' Puffer eine machbare Doppellektion (die separate
  Repetition der Feldformel des geraden Leiters entfällt, da bereits
  Vorwissen aus Magnetfelder). Falls trotzdem eng, erster Streichkandidat
  ist der historische Ampere-Exkurs (3'); als Nächstes die
  Feldlinien-Betrachtung (\glqq Gummifäden\grqq{}-Modell) als optionaler
  zweiter Erklärungsweg zur Leiterschaukel.
\end{itemize}

\end{document}
